% !TEX root = ../main.tex
\chapter{Discussion}
\label{ch:discussion}

Chapter~\ref{ch:results} reported the measurements taken on Arbitrum One. Here each research question and each objective of Chapter~\ref{ch:introduction} is held to the criterion stated there, and every interpretation points to the table of Chapter~\ref{ch:results} it rests on. Unless noted otherwise, same-window statements concern the matched cohort and holdout statements the spender cohort.

\dissertationSubheading[sec:findings-rq]{Findings by research question}

\dissertationSubsubheading[sub:finding-rq1]{RQ1: Reputation structure and rank divergence}

RQ1 asks whether EndorseRank and AWP order an identical sample of wallets in meaningfully different ways. The inter-method Kendall $\tau$ and its interval (Table~\ref{tab:alignment-family-ci}, last row) lie well inside the open interval between unrelated and identical orderings. The two scores are correlated, since both methods draw on overlapping wallet activity in one window, yet a large share of pairwise orderings disagree (Figure~\ref{fig:rank-scatter}).

Such divergence follows from how the two graphs are built. The endorsement graph records current spending authorization and has about one ninth as many edges as the transfer graph, which records time-decayed economic flow\cend{3}. Under the same damping convention, PageRank propagates influence along whatever edges it is given \cend{15}, and the way its stationary vector responds to the edge set has been worked out in detail \cend{23}. Different edge sets therefore place stationary mass differently. A wallet that receives large transfers but grants no allowances ranks high under AWP and modestly under EndorseRank; the reverse holds for a wallet that authorizes spenders without being a transfer hub. H1 is supported. In practice, one score cannot stand in for the other. An operator who adopts EndorseRank for its efficiency accepts a different ranking geometry and does not get a faster approximation of AWP.

\dissertationSubsubheading[sub:finding-rq2]{RQ2: Construct-specific alignment}

RQ2 asks which method aligns more strongly with contemporaneous allowance checks, and which with contemporaneous transfer checks. Tables~\ref{tab:alignment-family-ci} and~\ref{tab:tau-diff} show a different winner for each construct. AWP leads on the transfer family and EndorseRank on the allowance family, with between-method paired differences that exclude zero in opposite directions.

H2 has to be read with two caveats. One is a degree of methodological circularity in measuring each method against its own construct. Every proxy in a family is computed from the same event stream as the matching graph, so the construct validity seen under RQ2 is closer to validity in the sense of Cronbach and Meehl \cend{22} than to external validation. The other caveat is that no within-method contrast is significant: EndorseRank's allowance alignment is not significantly higher than its transfer alignment, and its agreement with in-approve degree is not significantly higher than its agreement with transfer in-degree (Table~\ref{tab:tau-diff}, contrasts~i and~v). The evidence thus supports the between-method statement alone, namely that transfer centrality does not reflect authorization and that endorsement centrality does not reflect transfer volume as well as AWP does. It gives no support to the claim that EndorseRank's own ranking leans more on allowances than on transfers. In this construct-specific, between-method form, H2 is supported.

\dissertationSubsubheading[sub:finding-rq3]{RQ3: Computational efficiency}

The question under RQ3 is whether EndorseRank is computationally more efficient than AWP on the same matched sample. On wall-clock time, memory and iteration count, Table~\ref{tab:benchmark-runtime} shows that it is. Edge counts are reported alongside the runtimes so that the ratio can be attributed to them: the endorsement graph is roughly one-ninth the size of the transfer graph, and each iteration costs time linear in $|E|$ \cend{23}. The graph is smaller because EndorseRank stores authorization as a snapshot of the latest allowance, where AWP keeps an accumulated, decayed history of transfers\cend{3}. Both methods run the same computation.

Table~\ref{tab:benchmark-scaling} qualifies this result more than it extends it. Across the different staged node counts that were tested, the ratio of wall-clock times stays within a narrow range. Each stage, though, adds address nodes without adding any extracted edges, so the stages do not test how either method scales once the edge population grows. Changing the damping or the top-token selection leaves the alignment pattern as it was but does shift the runtime ratio (Tables~\ref{tab:robustness-damping} and~\ref{tab:robustness-tokens}). H3 is supported on the matched cohort and within the limited experiments on node count.

\dissertationSubsubheading[sub:finding-rq4]{RQ4: Temporal holdout of future new approvals}

RQ4 has two clauses. The first asks whether EndorseRank scores computed up to the end of February predict new owner--spender approvals on inbound spenders over the following three months better than AWP does. According to Table~\ref{tab:holdout-spenders}, EndorseRank predicts these approvals and AWP does not. EndorseRank's confidence interval excludes zero, and its contrast with AWP is the strongest between-method contrast anywhere in the thesis (Table~\ref{tab:holdout-diff}). Because every label post-dates the scored events, this is no same-window construct check; it is the best evidence that endorsement centrality contains information about spenders that transfer centrality lacks. On the flow label the ordering repeats with a smaller margin, and in this cohort EndorseRank also predicts new transfer senders better than AWP.

The second clause asks whether EndorseRank also beats the same-window allowance degree. It does not. On new approval pairs the spender's in-approve degree at $t_1$ reaches $\tau=0.711$ against EndorseRank's $\tau=0.479$; the paired difference is $-0.233$, and its confidence interval lies far below zero. AWP stands in the same relation to transfer in-degree on new senders. The number of new counterparties a wallet gains within a quarter is correlated with the number it already had, and the local degree records that information without the interference that propagating the count brings in from owner-side centrality and allowance value. Whatever EndorseRank captures out of window is therefore inherited from the degree it smooths; it does not supplement that degree. None of this undoes the first clause, since out of window EndorseRank still differs from AWP by construct. It does rule out the possibility that propagation adds predictive power beyond the degree for this label, and the same is true of the transfer-graph baseline.

In one respect RQ4 departs from the research proposal: we made no comparison against a trading-success reference. That comparison was dropped because the outcome-native reference was derived from the same GMX events as the proxies it would have been compared with (Section~\ref{sub:gf-pr}). We compare instead against labels that lie outside the window but belong to the same construct, i.e.\ new approvals and new senders, and neither label contains any element of profit or repayment. What RQ4 shows is the temporal predictive validity of the endorsement construct relative to the transfer construct, with the raw degree as its upper bound. It does not address credit risk.

\dissertationSubsubheading[sub:finding-rq5]{RQ5: Coupling the two edge types}

RQ5 asks whether one walk over both layers can rank wallets better than either layer on its own. With the setting fixed before the run, it cannot. On the holdout, C-PR at $\lambda=0.5$ falls between EndorseRank and AWP on both labels; its paired difference to EndorseRank on new approval pairs is negative, and its paired difference to AWP on new senders includes zero (Table~\ref{tab:holdout-diff}). The sensitivity runs are ordered monotonically too, with $\tau$ rising on both labels as the allowance weight grows, which means that out of window the transfer layer adds no predictive power for spenders beyond what the allowance layer already supplies. We can also test C-PR against its secondary rule. In the same window (Section~\ref{sec:coupled-results}), C-PR tracks AWP to within a few thousandths on the transfer and Sybil-stability families and is indistinguishable from zero on the allowance family, so the secondary rule fails on its allowance part; the seeded ablation falls below the method it borrows from on every family. The denser layer dominates the mixture; $\lambda$ acts only at the minority of nodes with out-edges in both.

The ablation points to the cause. S-PR keeps the transfer walk but restarts it from the EndorseRank distribution; it reproduces AWP on the flow label and loses almost all of the endorsement signal on the approval label. When transfer edges govern propagation, a teleportation vector derived from endorsements changes little. The allowance edge produces its effect when it is walked on, not when it is used as a seed, and mixing it with transfer edges at the transition level dilutes it in proportion to the transfer weight. All of this concerns these two layers on this chain. We did not test a learned or node-dependent $\lambda$, nor a coupling that keeps the layers apart and combines their stationary vectors afterwards; both are left to future work (Chapter~\ref{ch:conclusion}).

We report this negative answer openly and count it as a contribution, for two reasons. Since the rule was pre-registered, the result cannot have been fitted to the data. It also puts a limit on an obvious next step for readers who take the crossed pattern of Table~\ref{tab:alignment-family-ci} as a reason to combine: on this evidence, a convex mixture of the two edge types at the transition level yields no score that beats the better single-layer method on either construct.

\dissertationSubheading[sec:objectives-achievement]{Achievement of the research objectives}

Chapter~\ref{ch:introduction} attached a measurement and a criterion to each of its five objectives. Table~\ref{tab:objectives-achievement} checks them against the evidence.

\begin{table}[htbp]
\centering
\caption{Research objectives, the criterion stated in Chapter~\ref{ch:introduction}, and the outcome.}
\label{tab:objectives-achievement}
\fitwidth{%
\begin{tabular}{p{0.06\linewidth}p{0.36\linewidth}p{0.42\linewidth}p{0.10\linewidth}}
\toprule
Obj. & Criterion & Evidence & Outcome \\
\midrule
O1 & 95\% interval of inter-method $\tau$ contains neither $0$ nor $1$ & Table~\ref{tab:alignment-family-ci}, last row: the interval lies strictly inside $(0,1)$ & Met \\
O2 & On its own construct, each method shows a positive $\Delta\tau$ over the other method, with an interval that excludes $0$ & Table~\ref{tab:tau-diff}, contrasts~ii (AWP on transfer) and~iii (EndorseRank on allowance) & Met \\
O3 & Runtime ratio given alongside the edge ratio, under a single fixed protocol & Tables~\ref{tab:benchmark-runtime}, \ref{tab:benchmark-scaling} and~\ref{tab:benchmark-scaling-incohort}: mean, SD, memory, iterations, $|V|$, $|E|$ per configuration & Met \\
O4 & EndorseRank's holdout interval excludes $0$ and its paired difference to AWP lies above $0$; the increment over $t_1$ in-approve degree is reported with its interval & Tables~\ref{tab:holdout-spenders} and~\ref{tab:holdout-diff}: first part met; the increment over degree is negative, and its interval lies below $0$ & Met (first clause); not met (second clause) \\
O5 & On the holdout, C-PR does no worse than EndorseRank on new approval pairs and better than AWP on new senders; in the same window, it lies within the leading single method's interval & Table~\ref{tab:holdout-diff}: both holdout parts fail; Table~\ref{tab:tau-diff-hybrid} covers the same window & Not met \\
\bottomrule
\end{tabular}
}
\end{table}

Three findings that lie outside the objectives also shape the conclusions. Because the within-method comparisons of Table~\ref{tab:tau-diff} were not statistically significant, less can be inferred about EndorseRank's class of methods (Section~\ref{sub:finding-rq2}). The expanded-pool experiment tested node count but not edge growth (Section~\ref{sub:finding-rq3}), and on the later labels of their own construct the raw $t_1$ degrees predict better than either PageRank (Section~\ref{sub:finding-rq4}). All three are reported as limits on the claims and are kept out of the objectives.

\dissertationSubheading[sec:theoretical-implications]{Theoretical implications}

The results bear on three areas of theory: propagation semantics on financial graphs, the construct validity of on-chain scores, and the economics of trust under pseudonymity.

\dissertationSubsubheading[sub:theory-propagation]{Propagation semantics beyond transfers}

PageRank was first formulated for hyperlinks, read as a citation network \cend{15}. Do and Do\cend{4} carried the metaphor over to Ethereum transactions, and Do, Do and Nguyen\cend{3} took it further to AWP, weighting transfer edges by value and by a logistic decay on age. EndorseRank applies the metaphor to ERC-20 allowances: an owner who lets a spender act for them in effect issues a directed endorsement of spending authority \cend{18}. As the RQ1 divergence shows, changing what edges mean changes the ranking space even with the same nodes. Graph-based reputation is thus a family of methods, each valid relative to the relationships its edges encode \cend{23}. Peer-to-peer and web trust propagation systems reached the same conclusion from the other side; in EigenTrust \cend{16} and TrustRank \cend{17}, the seed set and edge definitions largely decide the outcome. EndorseRank's new edge definition is native to a token standard, where the earlier ones belonged to a message-passing overlay or a hyperlink graph.

\dissertationSubsubheading[sub:theory-construct]{Multi-construct domain alignment}

The checks were designed to separate transfer centrality from allowance delegation and to test each method against both. Each method aligns more closely with its own construct and less with the other's. Classical construct-validity theory would predict this pattern for instruments that measure different constructs \cend{22}, and Messick's \cend{33} integrated view goes further, counting the consequences of using a score as part of its validity. Rank correlations \cend{20,21} put this logic into practice for ordinal reputation scores without requiring classification thresholds, while the paired bootstrap \cend{34} makes the pattern testable. We suggest this kind of formal comparison in preference to an informal reading of point estimates. The null within-method contrasts show that convergent evidence on a method's own construct does not amount to discriminant evidence within that method.

\dissertationSubsubheading[sub:theory-pseudonymity]{Reputation under cheap pseudonyms}

Identities are cheap to create on public ledgers \cend{5}, and under Ethereum's account model a new address needs no registration step at all \cend{36}. Any reputation system that rewards inbound transfers or allowances is therefore exposed to Sybil strategies, in which an attacker creates edges among addresses it controls \cend{11}. The empirical findings do not rule this out. They do show that the allowance and transfer graphs respond differently to the activity seen in the observational sample. In Section~\ref{sec:sybil-cost-model} we formalize the cost of attacking an allowance graph and argue that non-zero allowances to external spenders create opportunity and revocation costs that transfer wash trading does not incur.

\dissertationSubheading[sec:practical-implications]{Practical implications and deployment scenarios}

Protocols and infrastructure providers have to choose between transfer-hub detection and authorization-aware reputation. The measured trade-off supports choosing the score by scenario over recommending a single one.

\dissertationSubsubheading[sub:deploy-transfer]{Transfer-hub and flow monitoring}

A protocol whose priority is detecting high-flow wallets (for market-maker identification, routing analysis, or transfer-based compliance heuristics) has reason to choose AWP despite its cost. Given its transfer alignment (Table~\ref{tab:alignment-transfer}), AWP is the better fit when the operational target is economic flow centrality\cend{3}. Its runtime on the matched cohort (Table~\ref{tab:benchmark-runtime}) is adequate for batch analytics but less attractive for near real-time refresh.

\dissertationSubsubheading[sub:deploy-allowance]{Authorisation-aware and low-latency refresh}

Where explicit spending authorization matters most (token vaults, allowance-gated integrations, or risk engines that treat approvals as acts of trust), protocols have reason to prefer EndorseRank. Its allowance alignment (Table~\ref{tab:alignment-allowance}), its holdout result (Table~\ref{tab:holdout-spenders}), and its sparse graph with a short solve time (Table~\ref{tab:benchmark-runtime}) make it well suited to frequent recomputation. In return for authorization semantics, out-of-window evidence and speed, operators can accept weaker transfer alignment. Table~\ref{tab:holdout-diff} adds one caveat. If an operator's sole question is ``How many new approvals will a spender attract?'', the spender's current in-approve degree answers it better than any PageRank. The propagated score should be reserved for what it adds, which is resistance to the weight of a single owner and a ranking that reflects who the endorsing owners are.

\dissertationSubsubheading[sub:deploy-hybrid]{Evaluation design and combined use}

Whether to combine the two signals is a question about deployment and also about method, and the evidence answers the two differently. On the deployment side, an operator can run EndorseRank for continuous authorization-aware monitoring and AWP for periodic transfer-hub audits, at the cost of both solves added together. Folding the two edge types into a single walk finds no support. The coupled operator costs as much as AWP (Table~\ref{tab:benchmark-runtime}) and ranks no better than the appropriate single-layer method on either construct (Section~\ref{sub:finding-rq5}). An operator who wants both constructs should compute both scores and keep them apart. Because the pipeline is open, the analysis can be replicated across protocols on any ERC-20 chain that has allowance logs (Appendix~\ref{ch:appendix-eval}). Integration with live risk oracles is left for future work (Chapter~\ref{ch:conclusion}), although EndorseRank's speed could support a faster refresh cycle than the transfer-graph method allows at comparable cost.

\dissertationSubsubheading[sub:deploy-industry]{Industry adoption scenarios}

Either approach yields a cheap, standardized ordinal signal that others can recompute independently. The hypothesized sources of value are fees, utilization, monitoring cost and reduced reward leakage, not collateral. The scenarios below describe actors whose decisions the results address.

\begin{itemize}
\item DEX operators. A protocol that earns trading fees and refreshes its reputation scores frequently has to decide which wallets to prioritize for monitoring, fee tiers or feature access. EndorseRank fits this workflow because its score reflects authorization activity and can be recomputed in seconds without aggregating transfer history. Any gain in revenue would be realized through fee volume and lower monitoring cost, so fee and monitoring policies are the right place to apply the score.
\item Lending protocols. Lending markets want to put more of their capital to use without cutting collateral \cend{41}, and the stress events of 2020 showed what happens when risk parameters are updated late \cend{9}. Reputation can act as an extra signal for interest rates or fee tiers and for reviewing collateralized positions, though not as a replacement for collateral. For a lending protocol that already relies on ERC-20 approvals for deposits and borrowing (Section~\ref{sec:allowance-endorsement}), EndorseRank is the appropriate input. The scores are not tools for predicting default. Experiments with different collateral requirements are left to future work.
\item Wallet and security tooling. Users and integrators want to know whether a spender contract is widely authorized. A trust badge that shows the spender's EndorseRank, or its total incoming endorsement mass, makes explicit the delegation that allowance semantics encode \cend{18,19}. The badge conveys the position of a graph node under a defined set of edges; it is neither a security audit nor a guarantee against malicious contracts.
\item Liquidity providers and market makers. To control inventory and spreads, these participants need to find transfer hubs and see how flow is structured. AWP is the more suitable tool for that goal\cend{3}, and its benefit to them is lower screening cost, not higher profit and loss.
\item Airdrop and governance programmes. Sybil clusters routinely abuse token distributions and governance eligibility \cend{11,32}. Before rewards go out, programs can require diverse incoming endorsements, or they can apply the closed-cluster mass limit of Section~\ref{sec:sybil-cost-model} (Equation~\ref{eq:cluster-mass}) as an eligibility gate. Conditions of this kind limit the effect of closed-cluster inflation. They cannot stop open-cluster attacks that recruit honest incoming endorsements.
\item Analytics and compliance providers. Exchanges, funds and compliance teams need counterparty rankings that are auditable and that outsiders can recompute. The open source codebase allows EndorseRank and AWP scores to be published for defined windows and parameters; unlike proprietary machine learning models, these scores can be audited \cend{2}. Alongside any published ranking, providers should disclose its Sybil exposure and state that the scores are not profit predictions.
\end{itemize}

In none of these scenarios does reputation replace collateral; the standardized, auditable and low-latency ordinal signal lowers the cost of distinguishing between counterparties.

\dissertationSubsubheading[sub:deploy-multiples]{Standardised scores as shared infrastructure}

A standardized financial metric, a reputation score included, matters because it is public, comparable and cheap to recompute, so that several parties can base judgments on the same number and understand how it might be manipulated. Its predictive power matters less. EndorseRank and AWP meet the criteria of an infrastructure-level system for pseudonymous wallets, since both are computed from public ledgers with an open pipeline (Appendix~\ref{ch:appendix-eval}), can be refreshed quickly, and can be reproduced from the public ledger. Neither is meant as a forecast of profitability, and neither can be useful unless its computation is public and its vulnerabilities to manipulation are documented (Sections~\ref{sec:sybil-cost-model} and~\ref{sec:regulatory-ethics}).

\dissertationSubheading[sec:societal-impact]{Societal and practical impact}

In the DeFi market, vetting a counterparty is expensive, and the capacity to do it is unevenly spread. Wealthy actors may be able to pay for proprietary risk tools that analyze whole transfer histories, while retail participants work from anecdotal evidence or skip vetting altogether. A tool that uses the permissions users have already granted, and that can be updated for a fraction of the cost of re-analyzing the whole transfer graph (Table~\ref{tab:benchmark-runtime}), lowers the barrier to counterparty vetting. Most token-draining exploits rely on an approval that grants access to funds (Section~\ref{sec:approval-security}). A spender-centered endorsement score therefore conveys information about a risk users face directly and that transfer volume cannot measure.

For protocol developers, the results make a case for replacing a single ``reputation'' construct with specific scores, each fitted to the decision it informs. Such a change will reduce the risk of misapplication, for example using a flow score to decide authorization. Since both metrics can be recalculated from public data through a transparent pipeline, users will also be able to audit their own exposure. The lower computational cost will bring continuous monitoring within reach of protocols that cannot afford to recalculate the transfer graph periodically at the same expense.

For society at large, this paper shows that a useful on-chain reputation metric can be built with full transparency. Its input events are public and its parameters versioned, and every table can be rebuilt from scratch out of those inputs. Where the decentralized credit scoring systems discussed by Packin and Lev-Aretz\cend{2} are opaque, this system supplies the kind of explanatory information that regulators suggest should be required when DeFi protocols are analyzed \cend{63}. It also brings a duty, described in Section~\ref{sec:regulatory-ethics}, to ensure that an ordinal, graph-based statistic is not mistaken for an identity, a credit history, or proof of personhood.

The experiments ran on one blockchain over a six month period, so the results speak to the endorsement construct and not to credit risk. If they are to produce social value, they have to be deployed within the limits of that construct. Any reputation metric that affects how value is allocated is open to exploitation. The next section explains why manipulating these metrics is costly; it does not prove that manipulation is impossible.

\dissertationSubheading[sec:sybil-cost-model]{Sybil attack cost model on allowance graphs}

A Sybil attack exploits cheap pseudonyms to create a false impression of trustworthiness \cend{11}. On a transfer graph, a malicious actor can inflate in-degree and in-flow artificially by moving assets among accounts it controls, paying for this in transaction fees and possible inventory loss \cend{29}. Allowance graphs expose a different attack surface. Each edge there records an authorization, not an actual transfer. In this section we formulate a simple cost model for inflating EndorseRank with fake allowances.

\dissertationSubsubheading[sub:sybil-setup]{Attack setup}

Let the attacker control a set of addresses $\mathcal{A}$, and write $\mathcal{H}$ for the honest wallets in the endorsement graph $G_{\text{approve}}=(V,E,w)$. An edge $(u,v)$ of weight $w_{uv}>0$ records that $u$ has authorized $v$ to spend $w_{uv}$ tokens, counted after decimal adjustment. With damping parameter $d\in(0,1)$, EndorseRank finds the PageRank vector $\mathbf{r}$ on $G_{\text{approve}}$ as the solution of \cend{15}:
\begin{equation}
\label{eq:pagerank-sybil}
\mathbf{r} = d\, P^{\top} \mathbf{r} + \frac{1-d}{|V|}\mathbf{1},
\end{equation}
where $P$ is the row-stochastic matrix obtained by normalizing outgoing allowance weights, and dangling nodes are handled in the standard way \cend{23}. The attacker's objective is to maximize the mean EndorseRank of a target set $T\subseteq\mathcal{A}$,
\begin{equation}
\label{eq:objective}
J(T) = \frac{1}{|T|}\sum_{i\in T} r_i,
\end{equation}
subject to a budget constraint on the cost of the attack, which is defined next.

\dissertationSubsubheading[sub:sybil-cost]{Cost components}

Creating an allowance edge $(u,v)$ of weight $w_{uv}$ costs the attacker at least the following.

\begin{enumerate}
\item Gas cost. Each \texttt{approve} or \texttt{permit} call consumes gas. With $c_{\mathrm{gas}}$ the expected fee per allowance update, $m$ updates cost
\begin{equation}
C_{\mathrm{gas}} = m\, c_{\mathrm{gas}}.
\end{equation}

\item Capital lock and opportunity cost. An open allowance to an external spender $v\notin\mathcal{A}$ lets $v$ move the tokens held by $u$. Even with $v\in\mathcal{A}$, the attacker may still need large allowances to particular contracts or routers in order to look legitimate. Write $\ell(w_{uv})$ for the opportunity cost of the capital tied up in an edge of weight $w_{uv}$ over the period one reputation refresh covers (forgone interest, for example). Then
\begin{equation}
C_{\mathrm{cap}} = \sum_{(u,v)\in E_{\mathcal{A}}} \ell(w_{uv}),
\end{equation}
where $E_{\mathcal{A}}$ collects the edges that the attacker has created.

\item Revocation and detection cost. Legitimate users revoke old allowances from time to time, and an attacker who does not revoke might be identifiable. Let $c_{\mathrm{rev}}$ be the expected cost of periodically revoking and re-approving so as to keep a desired set of edges in place, for $k$ reputation refresh periods:
\begin{equation}
C_{\mathrm{rev}} = k\, c_{\mathrm{rev}}.
\end{equation}
\end{enumerate}

The total cost of the attack is the sum of the three components,
\begin{equation}
\label{eq:total-cost}
C = C_{\mathrm{gas}} + C_{\mathrm{cap}} + C_{\mathrm{rev}}.
\end{equation}

\dissertationSubsubheading[sub:sybil-bound]{Influence and a cost-effectiveness bound}

By the usual sensitivity results for PageRank, an edge moves $\mathbf{r}$ in proportion to the stationary mass at its source node and to its normalized outgoing weight \cend{23}. For a simple bound, suppose the attacker adds edges only among its own nodes, which yields a closed Sybil network that receives no endorsements from honest nodes. How far such edges inside $\mathcal{A}$ can raise $J(T)$ is then capped by the teleportation mass:
\begin{equation}
\label{eq:closed-cluster}
\sum_{i\in\mathcal{A}} r_i \leq |\mathcal{A}|\cdot\frac{1-d}{|V|} + d\cdot \sum_{i\in\mathcal{A}} r_i,
\end{equation}
which rearranges to
\begin{equation}
\label{eq:cluster-mass}
\sum_{i\in\mathcal{A}} r_i \leq \frac{|\mathcal{A}|}{|V|}.
\end{equation}
A purely closed cluster of attacker accounts therefore captures at most its population share of the total rank mass. Beating the bound requires incoming endorsements from $\mathcal{H}$. Obtaining them generally means either stealing the private keys of legitimate wallets or persuading their owners to authorize the attacker's accounts, and in both cases the cost is dominated by $C_{\mathrm{cap}}$ and by social-engineering effort, not by transaction fees.

Consider next an attacker who buys inbound allowances from legitimate owners, for instance through an app that asks users for permission. One unit of honest-sourced weight $w_{h a}$ then has a marginal cost that can be written
\begin{equation}
\label{eq:marginal}
c_{\mathrm{in}}(w_{h a}) = c_{\mathrm{gas}} + \ell_{\mathrm{honest}}(w_{h a}) + \pi(w_{h a}),
\end{equation}
where $\ell_{\mathrm{honest}}$ is the payment to the honest owner and $\pi$ is a risk premium against the allowance being spent by a bad actor. The attack's cost-effectiveness is then
\begin{equation}
\label{eq:efficiency}
\eta = \frac{\Delta J(T)}{C},
\end{equation}
with $\Delta J(T)$ the gain in the target set's mean rank that the attack edges produce.

Equations~\eqref{eq:closed-cluster}--\eqref{eq:efficiency} show that EndorseRank is not Sybil-proof. They do, however, bound closed-cluster inflation by population share, and they show that inflation across clusters requires paying for honest-sourced allowances. Transfer wash trading has other costs (fees and inventory), though it need not expose spending permissions to outside accounts. The Sybil-robust stability proxies of Chapter~\ref{ch:results} mitigate some forms of observational manipulation, but they cannot replace an adversarial assessment, which we leave to future work.

\dissertationSubheading[sec:regulatory-ethics]{Regulatory and ethical considerations}

On-chain reputation metrics raise questions of transparency as well as equity. Earlier legal work has shown that decentralized finance markets remain bound by disclosure and consumer protection rules even though settlement happens on-chain \cend{63}. Packin and Lev-Aretz\cend{2} add a warning of their own: decentralized credit scoring may perpetuate black-box harms when models cannot be interpreted or contested, or when they are linked to identity databases without due process. Both EndorseRank and AWP are computed deterministically from on-chain events with open-source algorithms, so anyone can recalculate them using the open-source software (Appendix~\ref{ch:appendix-eval}). That openness partly answers the black-box criticism. It leaves the distributional concern untouched, because addresses with little transaction activity will score lower whatever their off-chain creditworthiness.

Ethically, these scores should never be used as proxies for identity, as creditworthiness measures in the traditional sense, or as evidence that an address belongs to a unique individual \cend{11}. They are rank-based indicators of an address's position in the defined graph topology. If a deployment grants financial services on the strength of score thresholds, it should clearly document the graph definitions behind the scores, the timeframe, and the limitations, above all the exposure to Sybil attacks. A regulatory framework centered on interpretability would prefer this kind of ranking to proprietary machine learning models like those Kotb reviews \cend{14}, as long as providers make no exaggerated claims about how well they can predict loan repayment rates \cend{2}.

Privacy on an open blockchain is limited by the ledger's very nature \cend{5}, and computing the reputation scores draws on nothing beyond public data. The main ethical question is therefore how the results are put to use. Further attempts to uncover the identity behind each address are discouraged, and the scores should not be treated as permanent, fingerprint-like identifiers. This study works with wallet addresses alone and makes no attempt to link them to real-world individuals (Section~\ref{sec:ethics}).

\dissertationSubheading[sec:threats-validity]{Threats to validity}

The findings reported here should be read against the limitations below, which are grouped under construct, internal, external and adversarial validity.

\dissertationSubsubheading[sub:threat-construct]{Construct validity}

\begin{itemize}
\item These ratings are not credit ratings: no default labels from uncollateralized lending enter the analysis.
\item Same-window checks rest on local allowance and transfer statistics. In the holdout, the labels are future new approvals and future new transfer senders.
\item Holdout construct. A future new approval is still an endorsement event, just as a future new sender is still a flow event. Both lie outside the score window, yet neither is an independent credit or trading construct. Being counts of arrivals, they also favor any score that tracks the current count of counterparties (Section~\ref{sub:finding-rq4}).
\item Own-construct agreement is partly mechanical. EndorseRank smooths spender in-approve degree over the whole graph, and AWP does the same for decayed transfer in-strength, so each method is expected by construction to agree with its own family. That is why Table~\ref{tab:tau-diff} reports within-method contrasts, and why the null result on those contrasts is treated as informative.
\end{itemize}

\dissertationSubsubheading[sub:threat-internal]{Internal validity}

\begin{itemize}
\item Shared observation window. Every method and proxy is computed over 1~December~2025 to 31~May~2026. Market regimes peculiar to that window (volatility, incentive programs, protocol changes) may push alignment up or down. The robustness checks on damping and top-token subgraphs lower sensitivity to hyperparameters and token concentration, but none of them varies the calendar window.
\item Cohort selection. The matched cohort admits only wallets with non-zero connectivity in both reputation subgraphs. Results may not carry over to addresses absent from those subgraphs, and the sample-definition sensitivity in Table~\ref{tab:robustness-sample-size} shows how strongly the coefficients react once that requirement is relaxed.
\item Edge-weight measurement. Token-decimal normalization, unit conversion and permit-versus-approve coverage can all shift $w_{uv}$. Errors in the edge weights then propagate into $\mathbf{r}$ and into the rank correlations \cend{23}.
\item Timing environment. All runtimes are wall-clock times from a single machine. The standard deviations in Table~\ref{tab:benchmark-runtime} bound noise within a session but not variation between machines, so the comparison is stated in ratios instead of seconds.
\end{itemize}

\dissertationSubsubheading[sub:threat-external]{External validity}

\begin{itemize}
\item Single chain and window. The evidence comes from Arbitrum One alone, over six months. Extending the findings to other rollups, Layer-1 chains or lending-centered venues would need cross-protocol validation, since DeFi venues differ in protocol design \cend{8} and in market structure \cend{1}.
\item Pool scaling versus alignment scaling. The expanded pool supports only a statement about how runtime grows with node count (Section~\ref{sub:runtime-scaling}). Same-window alignment remains anchored to the matched cohort, and the holdout to spenders. The efficiency gains should not be taken to preserve the same $\tau$ profile at arbitrary population sizes.
\end{itemize}

\dissertationSubsubheading[sub:threat-adversarial]{Adversarial validity}

\begin{itemize}
\item Sybil and wash trading. Rankings may be manipulated through adversarial allowance or transfer patterns \cend{11}. The Sybil-adjusted proxies only partly address this risk, and the cost model of Section~\ref{sec:sybil-cost-model} is an analytical argument, not an empirical red-team evaluation.
\item Strategic approval behavior. Once EndorseRank is live, rational actors may grant allowances strategically to farm rank. An observational $\tau$ estimated before deployment may overstate how well the ranking resists such behavior.
\end{itemize}

\dissertationSubheading[sec:limitations]{Limitations}

All results were obtained over one six-month period on Arbitrum One, for a matched wallet cohort and a spender holdout cohort; Chapter~\ref{ch:results} reports the user counts and graph sizes for both. Among EndorseRank's strengths are construct alignment with allowances, a ranking distinct from AWP's, computational efficiency owed to graph sparsity, and, on the holdout, alignment with future new approvers where AWP shows none. Its weaknesses include weaker construct alignment than AWP with transfers and with longitudinal stability, an allowance lead over its own transfer alignment that did not reach significance, and a holdout coefficient below that of the raw in-approve degree it smooths.

The coupled operator was tried at three fixed values of $\lambda$, with uniform teleportation and a per-node convex mixture at the transition level. Node-dependent or learned weights, mixtures of the stationary vectors instead of the transition rows, and other Markov chains went untested, so the negative answer to RQ5 applies only to that family of methods on this dataset. The edge weights could also be influenced by how token decimals are normalized, by currency conversions, and by permit versus approve coverage. Because the second tier extraction did not finish, the expanded-pool experiment measures growth in nodes but not in edges. The longer-running baseline runs (other PageRank variants and a graph constructed around lending activity) are archived for future comparison (Appendix~\ref{ch:appendix-archive}) and left out of the main comparison. No default labels for unsecured lending exist in the dataset. Even where the results are construct-aligned, their deployment remains restricted by the regulatory and ethical considerations of Section~\ref{sec:regulatory-ethics}.

These limitations narrow what can be concluded. For this cohort and time frame, EndorseRank is a viable, computationally efficient, allowance-specific social reputation method; it is not a credit scoring method, nor a Sybil-resistant identity mechanism, nor a substitute for protocol-level collateralization \cend{9}. Chapter~\ref{ch:conclusion} restates what has been shown, what that knowledge is worth and what it would take to extend these bounds.
